\section{IPM: exact pressure, Fourier, and affine-wave correspondences}
\label{sec:ipm-exact-bridge}

This section audits equations (4.1)--(4.6) of the incoming research report,
physical PDF pages 10--11, and proves the additional identities used here.
Its starting equation is retained with the report's coefficients and signs:
\begin{equation}\label{eq:ipm-report-darcy}
 \partial_t\rho+u\cdot\nabla\rho=F,\qquad
 u=-\nabla p-\rho e_2,\qquad \nabla\cdot u=0,
 \qquad e_2=(0,1).
\end{equation}
The spatial coordinates are $x=(x_1,x_2)$; neither coordinate, time, nor a
coefficient in this equation is rescaled without an explicit map below.

\subsection{The primary-source coefficient and the Fourier zero mode}

The primary manuscript of Alp\"oge, Buckmaster, and Coiculescu,
\emph{Extending the C\'ordoba--Mart\'inez-Zoroa IPM Blow-up to Uniformly
Space-Time Smooth Forcing}, specifies on physical page 3, equations (1)--(2),
the torus $\mathbb T^2_{2\pi}=(\mathbb R/2\pi\mathbb Z)^2$ and operator
\begin{equation}\label{eq:ipm-source-operator}
 \widehat{u_{\mathbb T}(\sigma)}(k)=2\pi m(k)\widehat\sigma(k)
 \quad(k\in\mathbb Z^2\setminus\{0\}),\qquad
 \widehat{u_{\mathbb T}(\sigma)}(0)=0,
 \qquad \partial_t\sigma+u_{\mathbb T}(\sigma)\cdot\nabla\sigma=G,
\end{equation}
where
\begin{equation}\label{eq:ipm-symbol}
 \widehat f(k)=\frac1{(2\pi)^2}\int_{[0,2\pi]^2}f(x)e^{-ik\cdot x}\,dx,
 \qquad
 m(k)=\left(\frac{k_1k_2}{k_1^2+k_2^2},
             -\frac{k_1^2}{k_1^2+k_2^2}\right).
\end{equation}
The factor $2\pi$ is explicit in the source and absent in report equation
(4.2). The following results give the exact relation. The source's
Theorem 2.1 states smooth forcing on $[0,1]\times\mathbb T^2_{2\pi}$,
odd mean-zero density and force, smoothness before time $1$, and divergence
of density and velocity gradients. This statement is source evidence, not
an independent certification of that manuscript's full iteration.
The retrieved file, its hash, and the inspection locations are recorded in
the accompanying source inventory; the primary URL is
\url{https://cims.nyu.edu/~tristanb/ipm.pdf}.

\begin{proposition}[Pressure elimination with the zero mode retained]
\label{prop:ipm-pressure}
For any $\rho\in C^\infty(\mathbb T^2_{2\pi})$, set
$\overline\rho=\widehat\rho(0)$ and define the mean-zero periodic pressure
$\mathcal P\rho$ and the periodic vector field $\mathcal M\rho$ by
\begin{align}
 \widehat{\mathcal P\rho}(k)&=
       \frac{i k_2}{k_1^2+k_2^2}\widehat\rho(k)\quad(k\ne0),
 &\widehat{\mathcal P\rho}(0)&=0,\label{eq:ipm-pressure-fourier}\\
 \widehat{\mathcal M\rho}(k)&=m(k)\widehat\rho(k)\quad(k\ne0),
 &\widehat{\mathcal M\rho}(0)&=-\overline\rho e_2.
 \label{eq:ipm-full-operator}
\end{align}
Then $\mathcal M\rho=-\nabla\mathcal P\rho-\rho e_2$ and
$\nabla\cdot\mathcal M\rho=0$. Every smooth periodic solution of the
last two equations in \eqref{eq:ipm-report-darcy} has this velocity and
has pressure $\mathcal P\rho+c$ for a spatial constant $c$.
If $\mathcal M_0$ has the same nonzero symbol and zero symbol at $k=0$, then
\begin{equation}\label{eq:ipm-mean-pressure-relation}
 \mathcal M\rho=\mathcal M_0\rho-\overline\rho e_2,
 \qquad
 \mathcal M_0\rho=-\nabla\mathcal P\rho-(\rho-\overline\rho)e_2.
\end{equation}
On the universal cover $\mathbb R^2$, the same field $\mathcal M_0\rho$
satisfies the unprojected Darcy law with pressure
$\mathcal P\rho-\overline\rho x_2$. This pressure is periodic exactly when
$\overline\rho=0$.
\end{proposition}
\begin{proof}
Taking divergence of $u=-\nabla p-\rho e_2$ gives
$\Delta p=-\partial_2\rho$. For $k\ne0$ its Fourier coefficients obey
$-(k_1^2+k_2^2)\widehat p(k)=-i k_2\widehat\rho(k)$, giving
\eqref{eq:ipm-pressure-fourier}. Hence
\[
 \widehat u(k)=-ik\frac{ik_2}{k_1^2+k_2^2}\widehat\rho(k)
                   -e_2\widehat\rho(k)
 =\left(\frac{k_1k_2}{k_1^2+k_2^2},
 \frac{k_2^2}{k_1^2+k_2^2}-1\right)\widehat\rho(k)
 =m(k)\widehat\rho(k).
\]
The zero Fourier coefficient of a periodic gradient is zero, so
$\widehat u(0)=-\overline\rho e_2$, with no freedom to set it to zero
unless $\overline\rho=0$. Direct multiplication gives
$k\cdot m(k)=(k_1^2k_2-k_1^2k_2)/(k_1^2+k_2^2)=0$.
Smoothness follows because smooth Fourier coefficients decay faster than
every polynomial, whereas all displayed multipliers grow at most
polynomially; this also justifies every termwise derivative. The difference
of two periodic pressures with the same density has every nonzero Fourier
coefficient zero, and is therefore constant. Subtracting the zero-mode
formulas proves \eqref{eq:ipm-mean-pressure-relation}.
Finally $\nabla(-\overline\rho x_2)=-\overline\rho e_2$ proves the
cover-space formula. Translation by $2\pi e_2$ changes that pressure by
$-2\pi\overline\rho$, proving the stated periodicity criterion.
\end{proof}

\begin{proposition}[Exact coefficient map without changing space or time]
\label{prop:ipm-coefficient-map}
On the same torus and time interval, define the source and report residuals
\[
 \mathcal N_{2\pi}(\sigma)=\partial_t\sigma+
                   2\pi\mathcal M_0\sigma\cdot\nabla\sigma,
 \qquad
 \mathcal N_1(\rho)=\partial_t\rho+\mathcal M\rho\cdot\nabla\rho.
\]
For every smooth spatially mean-zero $\sigma$ the following is an exact
operator identity:
\begin{equation}\label{eq:ipm-coefficient-intertwining}
 \mathcal N_1(2\pi\sigma)=2\pi\mathcal N_{2\pi}(\sigma).
\end{equation}
Adjoin to the source scalar equation its periodic potential
$q=\mathcal P\sigma+c_{\rm s}(t)$, where $c_{\rm s}$ is an arbitrary smooth
function. It gives the exact velocity relation
$u_{\mathbb T}(\sigma)=2\pi(-\nabla q-\sigma e_2)$.
Consequently the invertible map of these full forced systems is
\begin{equation}\label{eq:ipm-source-report-map}
 \rho=2\pi\sigma,\quad F=2\pi G,\quad
 u=u_{\mathbb T}(\sigma)=2\pi\mathcal M_0\sigma,
 \quad p=2\pi q=2\pi\mathcal P\sigma+2\pi c_{\rm s}(t).
\end{equation}
It sends source initial density $\sigma^{\rm in}$ to
$\rho^{\rm in}=2\pi\sigma^{\rm in}$; its inverse divides density, force,
and pressure by $2\pi$ (recovering $q$) and leaves velocity, space, and time fixed.
It multiplies every density or force derivative and every homogeneous norm
of such a derivative by $2\pi$. Velocity and its derivatives are unchanged.
\end{proposition}
\begin{proof}
Linearity and the mean-zero condition give
$\mathcal M(2\pi\sigma)=2\pi\mathcal M_0\sigma$ and
$\mathcal P(2\pi\sigma)=2\pi\mathcal P\sigma$. Therefore
\[
 \partial_t(2\pi\sigma)+
 (2\pi\mathcal M_0\sigma)\cdot\nabla(2\pi\sigma)
 =2\pi[\partial_t\sigma+
        2\pi\mathcal M_0\sigma\cdot\nabla\sigma].
\]
Proposition~\ref{prop:ipm-pressure} proves both Darcy's law and
incompressibility for the pressure in \eqref{eq:ipm-source-report-map};
the gradient of $2\pi c_{\rm s}(t)$ is zero. Thus every spatially constant
pressure function is retained by the map and its inverse.
All inverse formulas follow by dividing by the nonzero constant $2\pi$.
Differentiation and absolute homogeneity of norms prove the last assertion.
Integrating the source transport equation on the torus gives
$\partial_t\overline\sigma=\overline G$, since the advective integral is
$\int\nabla\cdot(\sigma u_{\mathbb T}(\sigma))=0$; thus the source's
mean-zero forcing and initial density preserve the required subspace.
\end{proof}

There is also an exact treatment of the mean mode. Let $b(t)$ be any smooth
real function, let $B(t)=\int_{t_0}^t b(s)\,ds$, and write
$y=x+B(t)e_2$. For a smooth mean-zero source density $\sigma$ and an arbitrary
smooth real function $c(t)$ define
\begin{align}\label{eq:ipm-mean-dynamical-map}
 \rho(x,t)&=2\pi\sigma(y,t)+b(t),&
 u(x,t)&=u_{\mathbb T}(\sigma)(y,t)-b(t)e_2,\\
 p(x,t)&=2\pi\mathcal P\sigma(y,t)+c(t),&
 F(x,t)&=2\pi G(y,t)+b'(t).\notag
\end{align}
These are periodic in $x$ and solve \eqref{eq:ipm-report-darcy} whenever the
source residual is $G$. Indeed Darcy's law is
$-2\pi\nabla\mathcal P\sigma-(2\pi\sigma+b)e_2
=u_{\mathbb T}(\sigma)-be_2$ at $y$. Incompressibility is preserved by
translation. The transport residual is exactly
\[
 2\pi\sigma_t+2\pi B'\partial_2\sigma+b'
 +(u_{\mathbb T}(\sigma)-be_2)\cdot2\pi\nabla\sigma
 =2\pi G+b',
\]
because $B'=b$. Conversely, any smooth periodic report solution is put in
this form by $b=\overline\rho$, $c=\overline p$, $B(t)=\int_{t_0}^t b$, and
$\sigma(y,t)=[\rho(y-B(t)e_2,t)-b(t)]/(2\pi)$; the preceding equalities
read backwards prove its source equation. This retains the mean density,
the mean velocity, and their exact dynamical relation.

\subsection{Every mean-zero periodic phase profile}

\begin{proposition}[Wave, pressure, cancellation, and mean correction]
\label{prop:ipm-wave}
Let $k=(k_1,k_2)\in\mathbb Z^2\setminus\{0\}$, let
$h\in C^\infty(\mathbb R/2\pi\mathbb Z;\mathbb R)$, and suppose
$\int_0^{2\pi}h(s)\,ds=0$. Put $\kappa=k_1^2+k_2^2$ and let $H$ be the
unique mean-zero $2\pi$-periodic function with $H'=h$. For every real $a$,
\begin{align}\label{eq:ipm-periodic-wave}
 \vartheta(x)&=a h(k\cdot x),&
 v(x)&=a m(k)h(k\cdot x),&
 \pi(x)&=-\frac{a k_2}{\kappa}H(k\cdot x)
\end{align}
are smooth periodic fields satisfying
$v=\mathcal M\vartheta=\mathcal M_0\vartheta
=-\nabla\pi-\vartheta e_2$,
$\nabla\cdot v=0$, and $v\cdot\nabla\vartheta=0$.
For the primary source, velocity and pressure are respectively $2\pi v$
and $2\pi\pi$ at the unchanged source density $\vartheta$.
If the mean-zero hypothesis on $h$ is dropped and
$\overline h=(2\pi)^{-1}\int_0^{2\pi}h$, the report operator instead gives
\begin{equation}\label{eq:ipm-nonzero-mean-wave}
 \mathcal M[a h(k\cdot x)]
 =a m(k)(h(k\cdot x)-\overline h)-a\overline h e_2,
 \qquad
 \mathcal M[a h]\cdot\nabla(a h)
 =-a^2\overline h k_2h'(k\cdot x).
\end{equation}
\end{proposition}
\begin{proof}
The function $H_*(s)=\int_0^s h(r)\,dr$ is periodic because the integral
over a period is zero. Subtracting its mean gives $H$, and the difference
of two such primitives is a constant of mean zero. Write
$h(s)=\sum_{n\in\mathbb Z\setminus\{0\}}h_ne^{ins}$.
For each integer $N\ge0$, integration by parts $N$ times bounds
$|h_n|\le C_N|n|^{-N}$, so the series and all derivatives converge
absolutely. The frequency $nk$ is nonzero, and direct substitution gives
\[
 m(nk)=\left(\frac{n^2k_1k_2}{n^2(k_1^2+k_2^2)},
 -\frac{n^2k_1^2}{n^2(k_1^2+k_2^2)}\right)=m(k).
\]
Distinct $n$ give distinct frequencies $nk$. Thus the spatial mean is zero
and the multiplier acts on every harmonic by the same vector $m(k)$,
including negative $n$. Applying the pressure formula harmonic by harmonic,
or differentiating $\pi$, gives
\[
 -\nabla\pi-\vartheta e_2
 =a\left(\frac{k_1k_2}{\kappa},\frac{k_2^2}{\kappa}-1\right)h(k\cdot x)
 =a m(k)h(k\cdot x).
\]
Finally $\nabla\cdot v=a(m(k)\cdot k)h'=0$ and
$v\cdot\nabla\vartheta=a^2(m(k)\cdot k)hh'=0$.
The source assertion follows by multiplying velocity and pressure by the
explicit source coefficient. For nonzero mean, apply the proved formulas
to $h-\overline h$, retain the Darcy zero-mode term
$-a\overline h e_2$, and take its dot product with $ak h'$.
\end{proof}

The same direct pressure proof of \eqref{eq:ipm-periodic-wave} holds on
$\mathbb R^2$ for every nonzero real $k$, without imposing an integer lattice.
On a rectangular torus with periods $L_1,L_2>0$ it holds whenever
$k_jL_j\in2\pi\mathbb Z$ for $j=1,2$. These are explicit descent conditions,
not changes of coordinates. A profile with a smaller fundamental period
may also descend for additional $k$; no necessity claim is made for an
arbitrary non-faithful profile.

\subsection{The frozen growth rate is realized by a full affine solution}

\begin{theorem}[Exact unforced affine IPM solution and flow]
\label{thm:ipm-affine-solution}
Fix $A,a_0,t_0\in\mathbb R$, $k\in\mathbb R^2\setminus\{0\}$, and a real
smooth $2\pi$-periodic mean-zero profile $h$ with the primitive $H$ above.
Retain $\kappa=k_1^2+k_2^2$ and define
\begin{equation}\label{eq:ipm-affine-amplitude}
 \lambda=\frac{A k_1^2}{k_1^2+k_2^2},\qquad
 a(t)=a_0\exp\left(\frac{A k_1^2}{k_1^2+k_2^2}(t-t_0)\right).
\end{equation}
On the entire domain $\mathbb R^2\times\mathbb R$ set
\begin{align}\label{eq:ipm-affine-fields}
 \rho(x,t)&=A x_2+a(t)h(k_1x_1+k_2x_2),\\
 u(x,t)&=a(t)\left(\frac{k_1k_2}{k_1^2+k_2^2},
                  -\frac{k_1^2}{k_1^2+k_2^2}\right)h(k_1x_1+k_2x_2),\notag\\
 p(x,t)&=-\frac A2 x_2^2
         -\frac{a(t)k_2}{k_1^2+k_2^2}H(k_1x_1+k_2x_2).
 \notag
\end{align}
These fields solve the full system \eqref{eq:ipm-report-darcy} with $F=0$,
and are smooth at every finite time. The exact Lagrangian map and its inverse
are
\begin{align}\label{eq:ipm-affine-flow}
 \Phi_t(y)&=y+m(k)h(k\cdot y)I(t),&
 \Phi_t^{-1}(x)&=x-m(k)h(k\cdot x)I(t),&
 I(t)&=\int_{t_0}^t a(s)\,ds.
\end{align}
They preserve area and satisfy
$\rho(\Phi_t(y),t)=A y_2+a_0h(k\cdot y)$.
\end{theorem}
\begin{proof}
No pressure multiplier is applied to the affine function $Ax_2$ at
frequency zero: its pressure is the explicitly displayed polynomial.
Direct differentiation gives
\[
 -\nabla\left(-\frac A2x_2^2\right)-(Ax_2)e_2=0,
 \qquad
 -\nabla\left(-\frac{a(t)k_2}{\kappa}H(k\cdot x)\right)
       -a(t)h(k\cdot x)e_2=a(t)m(k)h(k\cdot x).
\]
Adding these identities proves Darcy's law. Write $s=k\cdot x$.
All terms of the divergence and transport equations are
\begin{align*}
 \nabla\cdot u&=a(t)h'(s)(k_1m_1(k)+k_2m_2(k))=0,\\
 \partial_t\rho&=a'(t)h(s),\\
 \nabla\rho&=Ae_2+a(t)k h'(s),\\
 u\cdot\nabla\rho
 &=Aa(t)m_2(k)h(s)+a(t)^2(m(k)\cdot k)h(s)h'(s)\\
 &=-\frac{A k_1^2}{k_1^2+k_2^2}a(t)h(s)+0.
\end{align*}
Since $a'=\lambda a$, their sum is identically zero.
For a spatial multi-index $\alpha=(\alpha_1,\alpha_2)$ and any integer
$j\ge0$, the wave part has derivative
\[
 \partial_t^j\partial_x^\alpha[a(t)h(k\cdot x)]
 =a_0\lambda^j e^{\lambda(t-t_0)}
     k_1^{\alpha_1}k_2^{\alpha_2}h^{(|\alpha|)}(k\cdot x).
\]
Together with the derivatives of the retained polynomial $Ax_2$ and
$-Ax_2^2/2$, this proves smoothness and boundedness of every derivative on
every compact subset of $\mathbb R^2\times\mathbb R$; wave derivatives are
also bounded on all space for every compact time interval.

Because $k\cdot m(k)=0$, $k\cdot\Phi_t(y)=k\cdot y$.
Differentiating in time yields
$\partial_t\Phi_t(y)=a(t)m(k)h(k\cdot y)=u(\Phi_t(y),t)$.
The same phase identity verifies both compositions of the displayed inverse.
The derivative matrix is
$D_y\Phi_t=\mathrm{Id}+I(t)h'(k\cdot y)m(k)\otimes k$.
Its two-dimensional determinant, expanded directly, is
\[
 (1+Ih'm_1k_1)(1+Ih'm_2k_2)-(Ih'm_1k_2)(Ih'm_2k_1)
 =1+Ih'(m_1k_1+m_2k_2)=1.
\]
Finally
\[
 \rho(\Phi_t(y),t)=Ay_2+
       [A m_2(k)I(t)+a(t)]h(k\cdot y).
\]
The derivative of the bracket is $A m_2(k)a(t)+a'(t)=0$ and its value at
$t=t_0$ is $a_0$, proving the transported-density identity. This argument
does not divide by $A$, $a_0$, or $k_1$, so the stationary degenerate cases
are included exactly.
\end{proof}

For $A>0$ and $k_1\ne0$ the report's rate
$-m(k)\cdot(Ae_2)=A k_1^2/(k_1^2+k_2^2)$ is therefore the exact growth rate
of this nonlinear solution. Its arbitrary-amplitude wave has identically
zero self-advection. The statement concerns an exponential in time, so it
establishes no finite-time singularity.

For completeness, retaining the primary coefficient in the cover-space
Darcy law $u=-\nabla p-2\pi\sigma e_2$ gives the exact fields
\begin{align}\label{eq:ipm-source-affine}
 \sigma&=A_{\rm s}x_2+a_{\rm s}(t)h(k\cdot x),\qquad
 u=2\pi a_{\rm s}(t)m(k)h(k\cdot x),\\
 p&=-\pi A_{\rm s}x_2^2-
       \frac{2\pi a_{\rm s}(t)k_2}{\kappa}H(k\cdot x),\qquad
 a_{\rm s}'=\frac{2\pi A_{\rm s}k_1^2}{\kappa}a_{\rm s}.
 \notag
\end{align}
Indeed the substitutions $A=2\pi A_{\rm s}$ and $a=2\pi a_{\rm s}$ in
\eqref{eq:ipm-affine-fields} give exactly these velocities and pressures,
and divide the scalar equation by $2\pi$. Thus the apparent rates have the
same value under the proved density map; no factor has been discarded.

\subsection{The exact domain and energy limitations}

If $A\ne0$, $\rho$ in \eqref{eq:ipm-affine-fields} cannot be a periodic
function on a full two-dimensional lattice. To prove this, suppose a lattice
period $\ell$ has $\ell_2\ne0$. Periodicity under every integer multiple
$n\ell$ would give
\[
 0=An\ell_2+a(t)[h(k\cdot x+n k\cdot\ell)-h(k\cdot x)].
\]
The second term has absolute value at most $2|a(t)|\|h\|_\infty$ whereas
$|An\ell_2|\to\infty$, a contradiction. A full two-dimensional lattice
contains such an $\ell$. When $A=0$ and $k\in\mathbb Z^2\setminus\{0\}$,
the formulas do descend to $\mathbb T^2_{2\pi}$ and reduce to the exact
stationary wave of Proposition~\ref{prop:ipm-wave}.

If $a_0\ne0$, $h\not\equiv0$, and $k_1\ne0$, the velocity has infinite
kinetic energy on $\mathbb R^2$ at every finite time. Here is the exact
integral. Set $J(k_1,k_2)=(-k_2,k_1)$ and use the invertible coordinates
$x=(s k+rJk)/\kappa$, so $k\cdot x=s$ and $dx=\kappa^{-1}ds\,dr$.
Moreover
\[
 |m(k)|^2=\frac{k_1^2k_2^2+k_1^4}{\kappa^2}
          =\frac{k_1^2}{\kappa}.
\]
For any $R>0$ the integral over $0\le s\le2\pi$, $-R\le r\le R$ is
\[
 \frac{2R}{\kappa}\,|a(t)|^2\frac{k_1^2}{\kappa}
                  \int_0^{2\pi}|h(s)|^2\,ds.
\]
All factors except $R$ are strictly positive, so letting $R\to\infty$
proves the assertion. If $k_1=0$, $m(k)=0$ and the velocity is zero; that
degenerate case is retained rather than included in the infinite-energy
claim. The affine scalar background also does not lie in $L^2$ when
$A\ne0$: if $k_1\ne0$ it grows linearly in $r$ in these same coordinates;
if $k_1=0$ it is a nonzero function of $x_2$ independent of $x_1$.
For the latter assertion, a nonzero affine function cannot be canceled
globally by the bounded periodic wave. The exact localized residual below
is the next map toward a finite-energy field; a failure of direct periodic
descent is only that descent obstruction.
